\documentclass[10pt]{article}

\usepackage[utf8]{inputenc}

\usepackage[T1]{fontenc}

\usepackage{graphicx}

\usepackage[export]{adjustbox}

\graphicspath{ {./images/} }

\usepackage{amsmath}

\usepackage{amsfonts}

\usepackage{amssymb}

\usepackage[version=4]{mhchem}

\usepackage{stmaryrd}

\usepackage{bbold}



\begin{document}

При изучении конкретных ф. я., заданных не в машинных терминах (напр., посредством формальных грамматик), часто возникает потребность нек-рым образом охарактеризовать сложность языка. Одним из наиболее распространенных путей в этом направлении является отыскание подходяпцего класса машин, распознающих рассматриваемые языки, и ошределение сложности языков через сложностные характеристики мапин. C другой стороны, изучение конкретного класса машин, как гравило, включает описание Ф. я., п. м. этого класса. Дальнейшее исследование Ф. я., п. м. затрагивает вопросы соотношения с известными классами языков, свойства замкнутости (относительно теоретико-множественных операций и т. п.), вопросы алгоритмического и сложностного характеров.



Лит.: [1] Языки и автоматы. Сб. переводов, М., 1975.



\includegraphics[max width=\textwidth, center]{2023_07_09_11ca9e6cd218c8317801g-1}

отношение между формальными системами, состоящее в том, что множества выражений, выводимых в әтих системах совпадают. Точнее, две формальные системы $S_{1}$ и $S_{2}$ әквивалентны тогда и только тогда, когда выполняются следующие условия: 1) всякая аксиома системы $S_{1}$ выводима в системе $S_{2} ; 2$ ) всякая аксиома системы $S_{2}$ выводима в системе $S_{1} ; 3$ ) если выражение $B$ непосредственно следует из выражений $A_{1}, \ldots, A_{n}$ в силу одного из правил вывода системы $S_{1}$ и выражөния $A_{1}, \ldots, A_{n}$ выводимы в системе $S_{2}$, то $B$ также выводимо в системе $S_{2}$; 4) аналогично 3) с заменой $S_{1}$ на $S_{2}$ и $S_{2}$ на $S_{1}$.



B. Е. Плиско



ФОРМУЛА - выражение борлализованного языка, предназначенное для записи суждения. Примеры точного определения понятия $Ф$. в различных формализованных языках см. в ст. Аксиоматическая теория множеств, Арифметика формальная, Іредикатов исчисление, Типов теория. В математич. практике Ф. наз. также осмысленные комбинации символов, несущие разнообразную смысловую нагрузку. Они могут быть как именными, так и высказывательными формами, определениями-сокращениями и пр.



ФОРТРАН (от ФОРмульный ТРАНслятор) - первоначально язык программирования задач вычислительной математики, разработанный в 1954-56 для машины ИБМ 704 и в версии Ф. II ставший первым в мире широко расгространенным алгоритмическим языком.



Позднее под Ф. стали понимать семейство языков программирования, восходящих к Ф. ІІ, прямо или косвенно претендующих на роль его преемника и сохраняющих название Ф. Большинство из них построено как расширение одного из двух стандартов (1962-64): Basic Fortran и Ф. IV и сходится в следуюпем.



Ф. предоставляет средства изображения, ввода и вывода арифметических, логических и текстовых данных. Арифметич. данные включают целые, действительные (обыкновенной и повышенной точности) и обычно такभе комплексные числа. Над арифметич. значениями определен обычный набор операций и отношений. Такие значения можно запоминать в скалярных переменных или әлементах однородных массивов (таблиц) того же типа.



Программа состоит из «ведущей программы» и набора подпрограмм (процедур), причем возможна их раздельная трансляция. Изменение текстуального порядка исполнения осуществляется посредством простых и вычисляөмых переходов по метке, условных операторов, циклов и вызовов процедур.



Подпрограммы передают друг другу информацию через свои параметры и общие блоки памяти (COMMON blocks). Так как рекурсивные вызовы запрещены, а границы массивов константны, часть действий по доступу $\kappa$ әлементам массивов может быть подготовлена при трансляции, что значительно ускоряет счет по готовой программе.



Можно сказать, что главная особенность Ф. заключена не в самом языке, а в той уникальной роли, к-рую он обеспечил себе в мировом программировании. Из «математического автокода» для ИБМ 704 он путем эволюции, напоминающей эволюцию естественных языков, превратился в самый распространенный язык программирования. Наличие хотя бы одной реализации Ф. стало практически обязательным для всякой ЭВМ общего назначения.



В әтой әволюции Ф. удержал ряд особенностей автокодов, важных для удобства программирования или использования программы, - таких, как составные константы (DATA) и раздельная трансляция. Они способствовали накоплению огромных библиотек ваписанных на Ф. подпрограмм по численным методам анализа, статистике, машинной графике, инженерным расчетам и т. Д.



В то же время Ф. сохранил весьма архаичный синтаксис. Группировку операторов (напр., в циклах) обеспечивают метки, что не только затемняет программу, но и ослабляет синтаксич. контроль. Поскольку описания идентификаторов необязательны, помехоустойчивость Ф. очень низка. Напр., если в типичном цикле



\[

\begin{aligned}

& D O 2 I=1,100 \\

& 2 A(I, J)=0

\end{aligned}

\]



пропустить любую запятую или заменить первую из них точкой, ошибка обнаружена не будет. Это снижает достоверность программ на Ф. и составляет важнеїший недостаток языка. Отчасти он преодолен в новых стандартах (напр., в Ф. 77).



Лuт.: [1] М а к - К р а к е н Д., Д о р н У., Численные методы и программирование на ФОРТР АНе, пер. с англ., 2 изд. М., 1977; 22 г р у н д Ф., Программирование на языке ФОРТран 77, пер с англ, м, 1982 С. Б. Покровский. ФОССА ПОВЕРХНОСТЬ - поверхность, несущая Фocca cemb.



ФОССА СЕТЬ - сопряженная геодезическая сеть. Поверхность, несупцая Ф. с., наз. поверхностью Фосса. На минимальной поверхности Ф. с. является изотропной сетью. Каждая Ф. с. на двумерной поверхности является главным основанием изгибания поверхности. Только геликоид несет бесконечное множество Ф. с. Вопрос о существовании поверхностей, несущих Ф. с., был поставлен А. Фоссом [1].



Jum. - [1] V o s s A., "Sitzungsber. math.-naturwiss. Kl. Bayerischen Akad. Wiss München», 1888, No 18, S 95-102; [2] Ф и н и к о в С. ІІ., Изгибание на главном основании,

М.- Л., 1937.



ФРАГМЕНА - ЛИНДЕЛЁФА ТЕОРЕМА - обобщение максимума модуля принципа аналитич. функций на случай функций, априори заданных как неограниченные; впервые в простейшей форме дано Э. Фрагменом и Ә. Линделефом [1]. Пусть $f(z)-$ регулярная аналитич. Функция комплексного переменного $z$ в области $D$ на плоскости $\mathbb{C}$ с границей $\Gamma$. Говорят, что $f(z)$ не превосходит по модулю числа $M$ в грангчной точке کЄГ, если



\[

\lim _{z \rightarrow \zeta} \sup _{z \in D}|f(z)| \leqslant M

\]



т. е. если для любого $\varepsilon>0$ найдется круг $\Lambda$ (зависящий от $\zeta$ и $\varepsilon)$ с центром $\zeta$ такой, что $|f(z)|<M+\varepsilon \operatorname{~іри~} z \in D \cap \Delta$. Основное содержание результатов Э. Фрагмена и Э. Линделёфа, в несколько модернизированной форме, заключается в следуюпцх двух положениях, последовательно расширяющих принцип максимума модуля.



I. Если регулярная аналитич. Функция $f(z)$ всюду на границе $\Gamma$ не превосходит по модулю числа $M$, то





\end{document}